\documentclass[11pt]{article}

\setlength{\parindent}{0pt}
\usepackage{xltxtra}
\usepackage{hyperref}
\hypersetup{hidelinks}
\usepackage{url}
\urlstyle{tt}
\usepackage{xcolor}
\definecolor{CVBlue}{RGB}{23,110,191}
\usepackage{calc}
\usepackage{graphicx}
\usepackage{tikz}
\usetikzlibrary{calc}
\usepackage{fontspec}
\usepackage{xeCJK}
\usepackage{enumitem}
\CJKsetecglue{} %% 取消中文与数字之间的间隙

%% 主文档字体设置（若本地无 fonts/ 目录，可改为系统字体）
\setmainfont[
    Path = fonts/Main/,
    Extension = .otf,
    BoldFont = texgyretermes-bold.otf,
]{texgyretermes-regular.otf}

\setCJKmainfont[
    Path = fonts/hansans/,
    Extension = .ttf,
    BoldFont = NotoSansSC-Bold.ttf,
]{NotoSansSC-Regular.otf}

\usepackage{relsize}
\usepackage{xspace}
\usepackage{fontawesome}

\usepackage[
    a4paper,
    left=1.2cm,
    right=1.2cm,
    top=1.5cm,
    bottom=1cm,
    nohead
]{geometry}

\renewcommand{\baselinestretch}{1.25}

\usepackage{titlesec}
\setlist{noitemsep}
\setlist[itemize]{topsep=0.25em, leftmargin=*}
\setlist[enumerate]{topsep=0.25em, leftmargin=*}

\newlength{\interProjectSpacing}
\setlength{\interProjectSpacing}{0.7em}
\newcommand{\projectsep}{\vspace{\interProjectSpacing}}

\newlength{\intraProjectTitleSep}
\setlength{\intraProjectTitleSep}{0.3em}
\newcommand{\titlebreak}{\\[\intraProjectTitleSep]}

\newlength{\intraProjectListTopSep}
\setlength{\intraProjectListTopSep}{0.15em}

\titleformat{\section}
{\large\bfseries\raggedright}
{}{0em}
{}
[{\color{CVBlue}\titlerule}]
\titlespacing*{\section}{0cm}{*1.2}{*0.9}

\begin{document}
\pagenumbering{gobble}

%%%% 照片（有 avatar.jpg 时取消注释）
% \begin{tikzpicture}[remember picture, overlay]
%   \node[anchor = north east] at ($(current page.north east)+(-2cm,-1.2cm)$)
%     {\includegraphics[height=3cm]{avatar.jpg}};
% \end{tikzpicture}

%%%% 校徽（有 logo 时取消注释）
% \begin{tikzpicture}[remember picture, overlay]
%   \node[anchor = north west] at ($(current page.north west)+(0.5cm,+1.0cm)$)
%     {\includegraphics[height=6cm]{ustc.png}};
% \end{tikzpicture}

\centerline{\LARGE\bfseries{单家豪}}

\vspace{0.3em}
\centerline{\normalsize{%
  \faPhone\ 13509655329 \quad
  \faEnvelopeO\ \href{mailto:jiahao020531@qq.com}{jiahao020531@qq.com} \quad
  \faMapMarker\ 广东深圳
}}
\centerline{\normalsize{生日：2002-05 \quad 政治面貌：共青团员}}

%% ==================== 教育背景 ====================
\section{\makebox[\widthof{\faGraduationCap}][c]{\color{CVBlue}\faGraduationCap}\ 教育背景}

\textbf{硕士：中国科学技术大学} \quad 前30\% \quad 软件工程 \hfill 2024.09 -- 2027.06
\begin{itemize}[nosep]
  \item 荣誉奖项：校级优秀共青团员，校级优秀学生干部
\end{itemize}

\vspace{0.4em}
\textbf{本科：哈尔滨工业大学（深圳）} \quad 前30\% \quad 计算机科学与技术 \hfill 2020.09 -- 2024.06
\begin{itemize}[nosep]
  \item 荣誉奖项：校级三等奖学金；第三届大湾区杯粤港澳金融数学建模竞赛本科组二等奖
\end{itemize}

%% ==================== 实习 / 项目经历 ====================
\section{\makebox[\widthof{\faUsers}][c]{\color{CVBlue}\faUsers}\ 实习经历}

\textbf{快影后端开发} \hfill 2025.09 -- 2026.04 \titlebreak
后端开发实习生
\begin{itemize}[nosep, topsep=\intraProjectListTopSep]
  \item \textbf{AI 智能素材标注与视频包装链路建设}
  \begin{itemize}[nosep]
    \item \textbf{AI 特征库建设}：负责素材智能标注链路设计与研发，对接底层 AI 工作流获取视频特征；基于 \textbf{Elasticsearch} 搭建正排索引支撑毫秒级检索，并同步落盘至 DB，打通 AI 特征上下游流转。
    \item \textbf{自动化视频包装}：扩展现有 AI 工作流系统，搭建以 AI 为核心的视频包装任务流；支持用户输入视频与风格配置后，自动检索 ES 素材特征并渲染生成 AI 视频效果，提升内容创作效率。
  \end{itemize}
  \item \textbf{千万级用户数据中心架构设计与优化}
  \begin{itemize}[nosep]
    \item \textbf{海量数据存储与同步}：针对 \textbf{5000 万}级用户底表，设计并落地 \textbf{MySQL 分表}方案；搭建定时流转任务，实现 Hive 到业务 DB 的高效增量同步，支撑下游高并发调用。
    \item \textbf{数据库性能调优}：引入「新用户空数据过滤」机制，拦截大量无效穿透查询，降低 DB 集群压力，保障数据中心高可用。
  \end{itemize}
  \item \textbf{商业化 DSP 广告归因系统对接}
  \begin{itemize}[nosep]
    \item \textbf{多维归因链路升级}：对接多家外部广告服务商（如巨量引擎），实现投放报表自动化拉取与转化数据回传；扩充下发维度与字段，提升多维归因覆盖率与准确性。
    \item \textbf{线上故障排查}：处理巨量引擎报表跨年拉取异常，定位并修复日期格式解析偏差，保障商业化结算数据准确性。
  \end{itemize}
  \item \textbf{搜索链路热词挖掘与线上稳定性治理}
  \begin{itemize}[nosep]
    \item \textbf{热搜词推荐与缓存加速}：搭建热搜词统计全链路，解析线上搜索日志并结合离线计算统计高频热词；引入空结果过滤，打通 Hive 至后端同步，并构建 \textbf{Local Cache} 实现极速查询。
    \item \textbf{慢查询排查与优化}：负责 Error 日志监控与链路追踪；针对慢查询 SQL 分析执行计划并建立联合索引，提升系统吞吐量。
  \end{itemize}
\end{itemize}

\projectsep

\textbf{百度智能云平台 · BMR 组件开发} \hfill 2025.06 -- 2025.08 \titlebreak
实习生\quad 百度智能云 - 通用数据部\\[0.2em]
参与 Baidu-MapReduce（BMR）大数据平台研发：BMR 基于 Spark / Hadoop 生态提供托管集群能力；\textbf{bmr-monitor} 负责状态监控与异常告警，\textbf{bmr-agent} 负责组件指标采集。
\begin{itemize}[nosep, topsep=\intraProjectListTopSep]
  \item \textbf{集群巡检数据采集}：实现巡检任务的周期调度与指标采集，沉淀可复用的巡检报告，支撑集群健康状态的持续观测。
  \item \textbf{智能巡检 RAG 应用}：基于 BCM 巡检报告指标编写 Prompt，搭建 \textbf{RAG} 问答链路，将结构化巡检结果转化为可交互的智能诊断能力。
  \item \textbf{性能压测支撑}：使用 Spark 提交 \textbf{TPC-DS} 标准基准任务，协助完成组件与集群数据处理性能评估。
\end{itemize}

%% ==================== 专业技能 ====================
\section{\makebox[\widthof{\faCogs}][c]{\color{CVBlue}\faCogs}\ 专业技能}
\begin{itemize}[nosep]
  \item \textbf{后端开发：} 熟悉 Java 核心与集合框架；熟练使用 Spring Boot、SpringMVC、MyBatis-Plus，了解 Spring IOC / AOP 原理。
  \item \textbf{数据存储：} 熟悉 MySQL 索引、事务与锁，具备分表设计、慢查询分析与联合索引优化经验；熟悉 Redis 缓存与常见高并发问题（穿透 / 击穿 / 雪崩）。
  \item \textbf{中间件与大数据：} 熟悉 Elasticsearch 检索与索引构建；了解 Hive 离线数仓同步及 Spark / Hadoop 生态基础。
  \item \textbf{计算机基础：} 熟悉 TCP/IP、HTTP/HTTPS；掌握操作系统进程通信、死锁等核心概念；具备线上日志排查与链路定位能力。
\end{itemize}
%% ==================== 社会实践与校园经历 ====================
\section{\makebox[\widthof{\faUsers}][c]{\color{CVBlue}\faUsers}\ 社会实践与校园经历}
\begin{itemize}[nosep]
  \item \textbf{社会实践}：积极参与社会志愿者活动，获得50个小时志愿时长。
  \item \textbf{校园经历}：研究生阶段：2024.09 -- 2025.09 担任24级软件工程网络与信息二班班级团支部的组织委员，2025年9月 -- 至今担任班级团支部副团支书；本科阶段：2020.09 -- 2024.06 担任2020级计算机科学与技术10班学习委员\textbf{校级优秀学生干部}。
\end{itemize}

\end{document}
